\documentclass[a4paper,11pt]{article}
\usepackage{exam-style}

\fancyhead[L]{\fontsize{9pt}{10pt}\selectfont\color{ctp-subtext0}341 农业综合知识三（信电）· 参考答案二}
\fancyhead[R]{\fontsize{9pt}{10pt}\selectfont\color{ctp-subtext0}信电学院部分}

\setlength{\parskip}{0.5ex}

\begin{document}
\examtitle

\parttitle{第一部分 \quad 程序设计基础知识（参考答案）}

\qtypename{一、填空题}{共5题，每题2分，共10分}

\q 第一空：\textbf{逻辑与}；第二空：\textbf{关系}（或 \code{==}、\code{!=} 等关系运算符）。关系运算符优先级高于逻辑与运算符。

\q \textbf{0（假）}。\code{a > b} 为假（3>4 不成立），逻辑与短路，结果为假（0）。

\q \textbf{0}。数组初始化部分元素，其余默认为 0。\code{a[1][2]} 是第二行第三列，未显式赋值。

\q 第一空：\textbf{5}；第二空：\textbf{6}。\code{strlen} 不计算终止符，\code{sizeof} 包含终止符。

\q \textbf{FILE *}（文件指针类型）。

\qtypename{二、简答题}{共5题，每题3分，共15分}

\q \code{break} 的作用是结束当前 \code{case} 的执行并跳出 \code{switch} 语句。如果省略 \code{break}，程序将继续执行后续 \code{case} 中的语句（即"穿透"效果），直到遇到 \code{break} 或 \code{switch} 结束。

\q \code{while} 循环：先判断条件，条件为真才执行循环体，可能一次都不执行。\code{do-while} 循环：先执行一次循环体，再判断条件，至少执行一次。

\q 数组名作为函数参数时，数组名退化为指向数组首元素的指针，即传递的是数组的地址，而不是整个数组的副本。因此在函数中对数组元素的修改会影响到原数组，且 \code{sizeof(arr)} 在函数内部得到的是指针的大小而非数组大小。

\q 数组是相同类型元素的集合，通过下标访问，所有元素类型相同。结构体由不同类型（异质）的成员组成，通过成员名访问。结构体支持成员的不同数据类型组合。

\q 文本方式：以字符为单位读写，自动进行换行符转换（\code{\textbackslash n} $\leftrightarrow$ \code{\textbackslash r\textbackslash n}）。二进制方式：按字节原样读写，不进行任何转换。

\qtypename{三、分析题}{共10题，每题2分，共20分}

\q \textbf{3 6 10}。case 1 穿透到 case 2 和 default：b=2+1=3；c=3+3=6；a+b+c=1+3+6=10。

\q \textbf{37}。累加 1$\sim$10 中不能被 3 整除的数：1+2+4+5+7+8+10=37。

\q \textbf{40 20}。\code{p = a + 3} 指向 40，\code{*(p - 2)} 指向 20。

\q \textbf{36}。累加所有元素：1+3+5+7+9+11=36。

\q \textbf{9 5}。通过指针传参实现了真正的交换，a=9, b=5。

\q
\texttt{\quad*}\\
\texttt{***}\\
\texttt{*****}\\
\texttt{*******}\\
（居中对齐的等腰三角形，4行）

\q \textbf{5 3}。通过加减运算交换变量值：a=5, b=3。

\q \textbf{-1}。\code{strcmp} 按字典序比较，"abc" < "abcd"，返回负数。

\q \textbf{5}。条件运算符 \code{? :} 取较大值 5。

\q \textbf{1 2 4 8 16}。每个元素是前一个的 2 倍。

\qtypename{四、程序设计题}{共2题，每题15分，共30分}

\pq 参考答案：
\begin{lstlisting}[language={[custom]C}]
#include <stdio.h>

int main() {
    int n, i, j, sum = 0;
    int a[10][10];

    printf("请输入矩阵阶数 n (n<=10): ");
    scanf("%d", &n);

    printf("请输入矩阵元素：\n");
    for (i = 0; i < n; i++)
        for (j = 0; j < n; j++)
            scanf("%d", &a[i][j]);

    for (i = 0; i < n; i++) {
        sum += a[i][i];                // 主对角线
        if (i != n - 1 - i)             // 避免交叉点重复加
            sum += a[i][n - 1 - i];    // 副对角线
    }

    printf("两条对角线元素之和为：%d\n", sum);
    return 0;
}
\end{lstlisting}

\pq 参考答案：
\begin{lstlisting}[language={[custom]C}]
#include <stdio.h>

void mystrcat(char *dest, const char *src) {
    while (*dest) dest++;        // 找到 dest 的末尾
    while (*src) {
        *dest = *src;
        dest++;
        src++;
    }
    *dest = '\0';                // 添加字符串结束符
}

int main() {
    char s1[100] = "Hello, ";
    char s2[] = "World!";
    mystrcat(s1, s2);
    printf("连接结果：%s\n", s1);
    return 0;
}
\end{lstlisting}

\parttitle{第二部分 \quad 数据库技术与应用（参考答案）}

\qtypename{一、填空题}{共5题，每题2分，共10分}

\q \textbf{树形}。层次模型使用树形结构表示实体及实体间的联系。

\q \textbf{笛卡尔积}。传统集合运算包括并、交、差和笛卡尔积。

\q \textbf{聚合}（或\code{COUNT、SUM、AVG}等聚合函数）。

\q \textbf{持久性（Durability）}。ACID 特性：原子性、一致性、隔离性、持久性。

\q \textbf{非平凡}。非平凡函数依赖指 Y 不是 X 的子集。

\qtypename{二、简答题}{共2题，每题5分，共10分}

\q 三级模式结构：外模式（用户视图）、概念模式（全局逻辑结构）、内模式（物理存储结构）。两级映射：外模式/概念模式映射（保证逻辑独立性）、概念模式/内模式映射（保证物理独立性）。映射使数据独立性得以实现，当某一层模式改变时只需修改映射关系，不影响其他层。

\q 数据独立性指应用程序不依赖于数据存储结构和存取方法的特性。物理独立性：当内模式（存储结构）改变时，只需修改概念模式/内模式映射，无需修改概念模式和外模式。逻辑独立性：当概念模式（全局逻辑结构）改变时，只需修改外模式/概念模式映射，无需修改应用程序。数据库通过三级模式和两级映射机制保障数据独立性。

\qtypename{三、操作题}{共8小题，共15分}

\q （2分）\code{SELECT Sno, Sname FROM S WHERE City = '北京';}

\q （2分）\code{SELECT DISTINCT Sno FROM SP WHERE Pno IN (SELECT Pno FROM P WHERE Color = '红色');}

\q （2分）\code{SELECT SUM(Qty) FROM SP WHERE Sno = 'S01';}

\q （2分）\code{SELECT Sno, Pno FROM SP WHERE Qty > 100;}

\q （2分）关系代数：\code{$\pi$_{Sno}(SP) $\div$ $\pi$_{Pno}(P)}（除法运算）

\q （2分）\code{SELECT Sno FROM SP WHERE Pno = 'P01' AND Sno IN (SELECT Sno FROM SP WHERE Pno = 'P02');} 或使用自连接。

\q （1分）\code{UPDATE S SET City = '广州' WHERE Sno = 'S03';}

\q （2分）
\begin{lstlisting}[language=SQL]
CREATE VIEW v_sp_info AS
SELECT Sname, Pname, Qty
FROM S, P, SP
WHERE S.Sno = SP.Sno AND P.Pno = SP.Pno;
\end{lstlisting}

\qtypename{四、分析题}{共1题，共10分}

\q （1）候选键：\textbf{(学号, 课程号)}。因为学号和课程号能函数决定所有属性。

\par（2）存在操作异常：
\begin{itemize}
  \item 插入异常：新学生尚未选课，无法插入学生信息（因为主键的一部分——课程号为空）。
  \item 删除异常：删除某学生的选课记录时，该学生的基本信息（姓名、系名等）也会被删除。
  \item 更新异常：如修改系主任信息，需要更新该系所有学生的记录。
\end{itemize}

\par（3）分解为 3NF：\\
学生表(\underline{学号}, 学生姓名, 系名)\\
课程表(\underline{课程号}, 学分)\\
系表(\underline{系名}, 系主任)\\
选课表(\underline{学号, 课程号}, 成绩)\\
（学号、课程号分别参照学生表和课程表，系名参照系表）

\qtypename{五、设计题}{共1题，共10分}

\q （1）E-R 图实体与联系：\\
实体：科室（科室编号，科室名称，科室位置）\\
实体：医生（工号，姓名，职称）\\
实体：患者（就诊卡号，姓名，性别，联系电话）\\
联系：属于（科室 $\to$ 医生，一对多）\\
联系：挂号（患者 $\leftrightarrow$ 科室，多对多；属性：挂号时间，就诊状态）

\par（2）关系模式（下划线为主键，* 为外键）：\\
科室表 \underline{科室编号}，科室名称，科室位置\\
医生表 \underline{工号}，姓名，职称，所属科室* REF 科室(科室编号)\\
患者表 \underline{就诊卡号}，姓名，性别，联系电话\\
挂号表 \underline{就诊卡号*，科室编号*}，挂号时间，就诊状态\\
（就诊卡号 REF 患者表，科室编号 REF 科室表）

\qtypename{六、应用题}{共2题，每题10分，共20分}

\q
\par （1）
\begin{lstlisting}[language=SQL]
SELECT city, COUNT(cid) AS cust_count,
       SUM(o.total) AS total_spent
FROM customers c
LEFT JOIN orders o ON c.cid = o.cid
GROUP BY city;
\end{lstlisting}
\par （2）
\begin{lstlisting}[language=SQL]
SELECT cname FROM customers
WHERE cid NOT IN (SELECT DISTINCT cid FROM orders);
\end{lstlisting}
\par （3）
\begin{lstlisting}[language=SQL]
DELIMITER $$
CREATE TRIGGER trg_update_total_spent
AFTER INSERT ON orders
FOR EACH ROW
BEGIN
    UPDATE customers
    SET total_spent = IFNULL(total_spent, 0) + NEW.total
    WHERE cid = NEW.cid;
END$$
DELIMITER ;
\end{lstlisting}

\q
\begin{lstlisting}[language=SQL]
DELIMITER $$
CREATE PROCEDURE update_stock(IN pid INT, IN qty INT)
BEGIN
    DECLARE cur_stock INT;
    SELECT stock INTO cur_stock FROM products WHERE pid = pid;
    IF cur_stock >= qty THEN
        UPDATE products SET stock = stock - qty WHERE pid = pid;
        SELECT '库存扣除成功' AS result;
    ELSE
        SELECT '库存不足' AS result;
    END IF;
END$$
DELIMITER ;
\end{lstlisting}

\end{document}
